\begin{tikzpicture}[
  node distance=10mm and 12mm,
  box/.style={draw, rounded corners, thick, align=center, minimum width=31mm, minimum height=10mm, fill=black!3},
  note/.style={draw, rounded corners, align=left, inner sep=4pt, fill=black!2},
  arrow/.style={-{Latex[length=2.5mm]}, thick, shorten <=2pt, shorten >=2pt}
]
\node[box] (access) {State access\\and scheduling};
\node[box, right=of access] (exec) {State-aware\\execution};
\node[box, right=of exec] (evo) {State evolution\\and reuse};
\path (access.north) ++(0,9mm) coordinate (feedbackL);
\path (evo.north) ++(0,9mm) coordinate (feedbackR);

\draw[arrow] (access) -- node[pos=0.46, above=5pt, align=center, fill=white, inner sep=1pt, font=\small] {conflicts, locality,\\queueing} (exec);
\draw[arrow] (exec) -- node[pos=0.54, above=5pt, align=center, fill=white, inner sep=1pt, font=\small] {cost, quality,\\resource usage} (evo);
\draw[arrow] (evo.north) -- (feedbackR) -- node[midway, above=2pt, align=center, fill=white, inner sep=1pt, font=\small] {future hotspots,\\memory pressure} (feedbackL) -- (access.north);

\node[note, below=12mm of exec, text width=107mm] (note1) {
\textbf{Propagation view.} Local mismatches rarely stay local. Hidden access contention distorts execution decisions; poorly modeled execution reshapes what can be updated or retained; aggressive updates feed back into later access pressure and reuse quality.
};

\draw[arrow] (exec.south) -- (note1.north);
\end{tikzpicture}
